%!TEX TS-program = lualatex
\documentclass[a0paper,portrait,25pt]{tikzposter}

% Поддержка кириллицы и шрифтов для LuaLaTeX
\usepackage{fontspec}
\usepackage{polyglossia}
\setmainlanguage{russian}
\setotherlanguage{english}

\setmainfont{DejaVu Sans}
\setsansfont{DejaVu Sans}
\setmonofont{DejaVu Sans Mono}
\newfontfamily\cyrillicfont{DejaVu Sans}
\newfontfamily\cyrillicfontsf{DejaVu Sans}
\newfontfamily\cyrillicfonttt{DejaVu Sans Mono}

\usepackage{amsmath,amssymb}
\usepackage{graphicx}
\usepackage{tikz}
\usetikzlibrary{shapes.geometric, arrows.meta, positioning}

% Базовая тема
\usetheme{Simple}
\usecolorstyle{Default}

% Принудительная настройка контрастных цветов (белый фон, синие акценты)
\colorlet{backgroundcolor}{white}
\colorlet{framecolor}{blue!70!black}
\colorlet{titlefgcolor}{blue!80!black}
\colorlet{blocktitlefgcolor}{blue!80!black}
\colorlet{blocktitlebgcolor}{white}
\colorlet{blockbodybgcolor}{white}
\colorlet{blockbodyfgcolor}{black}

% Метаданные с центрированием и фиксацией ширины
\title{\parbox{\linewidth}{\centering \textbf{Разработка масштабируемых систем\\визуализации данных}}}
\author{Тоц Леонид Александрович}
\institute{Группа ИВТ-2}
\usepackage{pgfplots}\pgfplotsset{compat=1.18}

\begin{document}

\maketitle

% Блок аннотации
\block{Аннотация (Abstract)}{
    \Large
    В работе исследуются алгоритмы и конвейеры автоматизированной вёрстки сложных аналитических данных.
    Рассмотрены подходы к представлению многоуровневых структурных графов, векторных диаграмм и формированию
    отчётных постеров формата A0 без потери типографического качества при масштабировании.
}

\begin{columns}
    % ================= ЛЕВАЯ КОЛОНКА =================
    \column{0.5}

    \block{1. Введение и постановка задачи}{
        \large
        При визуализации результатов научных исследований возникает задача компактного и наглядного размещения графики, математических выражений и программного кода.

        \vspace{0.8em}
        \textbf{Основные цели исследования:}
        \begin{itemize}
            \item Обеспечение воспроизводимости графиков и блок-схем на базе декларативного кода \LaTeX;
            \item Исключение артефактов масштабирования за счёт сквозного векторного рендеринга;
            \item Оптимизация компоновки блоков формата A0.
        \end{itemize}
    }

    \block{2. Архитектура и методы генерации}{
        \large
        Конвейер генерации визуализации состоит из трёх последовательных этапов:

        \vspace{1.2em}
        \begin{center}
        \begin{tikzpicture}[
            node distance = 1.2cm,
            box/.style = {
                rectangle,
                rounded corners=10pt,
                minimum width=6.2cm,
                minimum height=3.2cm,
                align=center,
                font=\large\bfseries,
                fill=blue!10,
                draw=blue!80!black,
                line width=2pt
            },
            arrow/.style = {-{Stealth[scale=1.3]}, line width=2.5pt, draw=blue!80!black}
        ]
            \node[box] (in) {Входные данные\\(CSV / JSON)};
            \node[box, fill=teal!15, draw=teal!80!black, right=of in] (proc) {Транслятор\\TikZ / PGF};
            \node[box, fill=orange!15, draw=orange!80!black, right=of proc] (out) {Векторный\\макет A0};

            \draw[arrow] (in) -- (proc);
            \draw[arrow] (proc) -- (out);
        \end{tikzpicture}
        \end{center}

        \vspace{1em}
        Математическая оптимизация размещения элементов графа:
        \begin{equation*}
            E(p) = \sum_{(u,v) \in \mathcal{E}} \|p_u - p_v\|^2 + \lambda \sum_{u \neq v} \frac{1}{\|p_u - p_v\|}
        \end{equation*}
    }

    \block{3. Практические результаты}{
        \large
        Применение декларативной вёрстки позволило сократить время подготовки финальной графики к публикации на 42\% по сравнению с классическими растровыми редакторами.
    }

    % ================= ПРАВАЯ КОЛОНКА =================
    \column{0.5}

    \block{4. Анализ производительности}{
        \large
        Оценка времени компиляции и нагрузки на память при различном количестве узлов графа:

        \vspace{0.8em}
        \begin{center}
            \begin{center}
                \begin{tikzpicture}
                    \begin{axis}[
                        ybar,
                        bar width=1.5cm,
                        width=0.75\linewidth,
                        height=14cm,
                        ylabel={\textbf{Время рендеринга (мс)}},
                        symbolic x coords={50 узлов, 200 узлов, 1000 узлов},
                        xtick=data,
                        nodes near coords,
                        nodes near coords align={vertical},
                        every node near coord/.append style={font=\large\bfseries},
                        ymin=0, ymax=600,
                        ymajorgrids=true,
                        grid style=dashed,
                        legend style={at={(0.5,1.05)}, anchor=south, legend columns=-1, font=\Large},
                        tick label style={font=\large},
                        label style={font=\Large\bfseries}
                    ]
                        \addplot[draw=blue!80!black, fill=blue!40] coordinates {(50 узлов, 45) (200 узлов, 180) (1000 узлов, 510)};
                        \addplot[draw=teal!80!black, fill=teal!40] coordinates {(50 узлов, 22) (200 узлов, 78) (1000 узлов, 240)};
                        \legend{Базовый метод, Предложенный метод}
                    \end{axis}
                \end{tikzpicture}
            \end{center}
        \end{center}

        \vspace{0.8em}
        \textbf{Ключевые выводы тестирования:}
        \begin{enumerate}
            \item Высокая точность позиционирования типографических гарнитур;
            \item Изоляция слоёв диаграмм исключает взаимное перекрытие блоков;
            \item Полная поддержка векторного стандарта при экспорте в PDF.
        \end{enumerate}
    }

    \block{5. Заключение}{
        \large
        Использование пакета \texttt{tikzposter} в связке с компилятором \textbf{LuaLaTeX} обеспечивает стабильную сборку сложных академических постеров с поддержкой русского языка. Шаблон пригоден для дипломных докладов и постерных конференций.
    }

    \block{6. Список литературы}{
        \normalsize
        \begin{enumerate}
            \item \textbf{Tantau T.} \textit{The TikZ and PGF Packages}. Manual for version 3.1.10, CTAN, 2023.
            \item \textbf{Richter C.} \textit{The tikzposter package: Guidelines and documentation}. CTAN, 2014.
            \item \textbf{Knuth D. E.} \textit{The TeXbook}. Addison-Wesley, 1986.
        \end{enumerate}
    }

\end{columns}

\end{document}